% ---------------------------------------------------------------------------
% Ablation of the v3 prior redesign.
%
% DRAFTING NOTE: this is deliberately long. Every number is recorded here with
% the epoch it was measured at so the section can be cut down later without
% having to re-derive anything. Source data:
%   runs/prior_h100_cpc_v3/validation/history.jsonl           (A1)
%   runs/prior_h100_cpc_v3_hurdle/validation/history.jsonl    (A2)
%   runs/prior_h100_cpc_v3_absolute/validation/history.jsonl  (A3)
%   data/processed/v2_ingestion_triplet/ing2022_s04_g010_sqrtfix/
%     v1_vs_v2/v1_vs_v2_selection.{json,md}                   (V1--V2 DA)
% ---------------------------------------------------------------------------

\section{Ablation: components of the prior redesign}
\label{sec:ablation-v3}

\subsection{Motivation}

The baseline prior (hereafter V1) reproduces daily rainfall amount acceptably but
rains far too often. Over 400 evaluation days it produces measurable rainfall on
$64.9\%$ of land cells against CHIRPS's $45.9\%$, a wet-frequency error of
$+0.19$. The error is concentrated where it matters most for assimilation:
because the BMD gauge network sits in the drier half of the domain (CHIRPS
averages $6.7$~mm~day$^{-1}$ at station locations against $12.3$~mm~day$^{-1}$
domain-wide), the resulting bias is $+5.88$~mm~day$^{-1}$ at stations against
$+2.65$~mm~day$^{-1}$ on the grid. A background that rains everywhere lightly
biases every increment the analysis subsequently makes.

Two mechanisms were hypothesised, and the ablation is designed to separate them.

\begin{enumerate}
  \item \textbf{Continuous density.} Rectified flow models a continuous
    density, whereas daily rainfall has an atom at exactly zero carrying
    ${\sim}54\%$ of its mass. A continuous model can only approach dry
    asymptotically, so dry cells decode to small positive rainfall. The proposed
    remedy is a \emph{hurdle head}: a second output channel predicting
    $P(\text{dry})$, trained with binary cross-entropy and imposed as a mask
    after sampling.
  \item \textbf{Moving target.} V1 predicts the residual
    $T(\mathrm{CHIRPS}) - T(\mathrm{CPC})$. Because the CPC base is wet over
    $84\%$ of the domain, ``dry'' is not a fixed value in residual space but a
    quantity that varies with the base. The proposed remedy is to predict the
    \emph{absolute} transformed field.
\end{enumerate}

\subsection{Design}

Four configurations were trained for 150 epochs each, identical in
architecture ($51.8$M parameters), optimiser, schedule, data
(\texttt{bd\_wide\_cpc.zarr}) and normalisation statistics
(\texttt{stats\_cpc.json}). They differ only as shown in
Table~\ref{tab:ablation-design}.

\begin{table}[htbp]
\centering
\caption{Ablation design. Each arm differs from the full redesign (A1) in
exactly one respect, so the contribution of each component is identifiable.}
\label{tab:ablation-design}
\begin{tabular}{llccc}
\hline
Arm & Description & Target & Hurdle head & \texttt{dry\_weight} \\
\hline
A1 & full redesign        & absolute & yes & 2.0 \\
A2 & hurdle head only     & residual & yes & 1.0 \\
A3 & absolute target only & absolute & no  & 1.0 \\
A4 & residual + full loss & residual & yes & 2.0 \\   % cond_dropout 0.0
V1 & baseline             & residual & no  & --- \\
\hline
\end{tabular}
\end{table}

The \texttt{dry\_weight} parameter upweights dry cells in the flow-matching
loss. It is set to $2.0$ only in A1, on the reasoning that an unweighted mean
over a field that is half dry still lets the wet half dominate the gradient,
since wet cells carry much larger residuals. As Section~\ref{sec:abl-findings}
shows, this parameter turns out to matter more than the hurdle head itself, and
the fact that it is confounded with the target change in A1 is a limitation of
this design.

\paragraph{Validation protocol.} Every five epochs, an eight-member ensemble is
drawn with EMA weights over 30 sampler steps at classifier-free guidance weight
$w=1$, on six held-out July days selected to span the domain-mean rainfall
distribution: 2019-07-20 (q05), 2020-07-17 (q25), 2020-07-31 (q50), 2020-07-10
(q75), 2019-07-26 (q95) and 2020-07-11 (q99). Wet-day frequency is evaluated at
a $1$~mm~day$^{-1}$ threshold. Observed wet fraction on these six days is
$0.696$; note this is a monsoon subset and is therefore much wetter than the
annual figure of $0.459$ against which V1's $+0.19$ error was originally
quoted. Errors, not absolute frequencies, are the comparable quantity.

\subsection{Results}

\begin{table}[htbp]
\centering
\caption{Ablation results at the best-CRPS epoch and at epoch 149. Wet-fraction
error is predicted minus observed ($0.696$ on the validation days); zero is the
target. Spatial correlation and bias are against CHIRPS, which is the training
target and therefore the appropriate reference during training. Bias in
mm~day$^{-1}$.}
\label{tab:ablation-results}
\begin{tabular}{llrrrrrr}
\hline
Arm & Epoch & CRPS & Wet frac. & Wet err. & Corr. & Bias & Spread \\
\hline
A1 & 24 (first)  & 19.08 & 0.629 & $-0.067$ & 0.075 & $-19.44$ & 4.53 \\
A1 & 39          & ---   & 0.698 & $+0.002$ & ---   & ---      & ---  \\
A1 & 134 (best)  & 14.49 & 0.920 & $+0.224$ & 0.140 & $-12.04$ & 14.35 \\
A1 & 149 (final) & 14.56 & 0.922 & $+0.226$ & 0.132 & $-12.24$ & 14.05 \\
\hline
A2 & 24 (first)  & 10.63 & 0.997 & $+0.301$ & 0.325 & $-1.44$  & 14.89 \\
A2 & 74 (best)   & 10.40 & 0.995 & $+0.299$ & 0.301 & $+5.69$  & 21.06 \\
A2 & 149 (final) & 10.97 & 0.994 & $+0.298$ & 0.286 & $+8.82$  & 23.35 \\
\hline
A3 & 9 (first)   & 18.96 & 0.994 & $+0.299$ & 0.174 & $-18.43$ & 2.24 \\
A3 & 144 (best)  & 13.37 & 0.965 & $+0.270$ & 0.119 & $-10.29$ & 14.63 \\
A3 & 149 (final) & 13.41 & 0.966 & $+0.270$ & 0.125 & $-10.49$ & 14.29 \\
\hline
A4 & 9 (first)   & 14.90 & 0.892 & $+0.196$ & 0.404 & $-13.69$ & 5.35 \\
A4 & 39 (best)   & 9.38  & 0.967 & $+0.272$ & 0.349 & $+1.58$  & 20.84 \\
A4 & 149 (final) & 10.83 & 0.991 & $+0.295$ & 0.293 & $+9.24$  & 27.11 \\
\hline
V1 & ---         & ---   & ---   & $+0.19$  & ---   & $+2.65$  & --- \\
\hline
\end{tabular}
\end{table}

A4 was run after the guidance probe of F6 and sets
$\texttt{cond\_dropout}=0.0$, matching V1; the other three use $0.1$.

\begin{table}[htbp]
\centering
\caption{Full validation trajectories, every fifteenth epoch. For each arm:
CRPS (mm~day$^{-1}$), wet-frequency error, ensemble spread (mm~day$^{-1}$) and
spatial correlation. A1 and A2 were restarted at epoch 20 and so have no epoch-9
entry. Recorded in full so the section can be shortened later without returning
to the source histories.}
\label{tab:ablation-trajectories}
\small
\begin{tabular}{r|rrrr|rrrr|rrrr|rrrr}
\hline
 & \multicolumn{4}{c|}{A1 abs+hurdle} & \multicolumn{4}{c|}{A2 res+hurdle}
 & \multicolumn{4}{c|}{A3 abs only} & \multicolumn{4}{c}{A4 res+full} \\
Ep & CRPS & werr & sprd & $r$ & CRPS & werr & sprd & $r$ & CRPS & werr & sprd & $r$ & CRPS & werr & sprd & $r$ \\
\hline
9 & --- & --- & --- & --- & --- & --- & --- & --- & 19.0 & +0.299 & 2.2 & 0.17 & 14.9 & +0.196 & 5.4 & 0.40 \\
24 & 19.1 & -0.067 & 4.5 & 0.08 & 10.6 & +0.301 & 14.9 & 0.33 & 16.7 & +0.278 & 5.8 & 0.16 & 10.5 & +0.252 & 15.0 & 0.35 \\
39 & 18.4 & +0.002 & 6.5 & 0.06 & 10.5 & +0.303 & 18.3 & 0.28 & 16.4 & +0.258 & 7.0 & 0.15 & 9.4 & +0.272 & 20.8 & 0.35 \\
54 & 17.8 & +0.023 & 8.1 & 0.05 & 10.5 & +0.302 & 19.6 & 0.29 & 15.8 & +0.251 & 7.8 & 0.15 & 9.4 & +0.281 & 21.8 & 0.34 \\
69 & 16.8 & +0.120 & 8.7 & 0.12 & 10.4 & +0.301 & 20.3 & 0.30 & 15.3 & +0.256 & 8.8 & 0.14 & 9.9 & +0.286 & 23.0 & 0.31 \\
84 & 16.1 & +0.148 & 9.9 & 0.12 & 10.6 & +0.297 & 21.1 & 0.29 & 14.7 & +0.264 & 10.0 & 0.16 & 10.1 & +0.290 & 24.5 & 0.32 \\
99 & 15.3 & +0.186 & 11.3 & 0.13 & 10.9 & +0.299 & 22.1 & 0.29 & 14.2 & +0.264 & 11.4 & 0.15 & 10.3 & +0.289 & 25.8 & 0.31 \\
114 & 14.9 & +0.199 & 13.2 & 0.13 & 11.1 & +0.299 & 22.9 & 0.29 & 13.6 & +0.268 & 13.2 & 0.15 & 10.5 & +0.291 & 26.6 & 0.31 \\
129 & 14.5 & +0.221 & 14.3 & 0.14 & 11.2 & +0.299 & 23.4 & 0.29 & 13.5 & +0.268 & 14.4 & 0.12 & 10.7 & +0.293 & 26.9 & 0.30 \\
144 & 14.5 & +0.228 & 14.1 & 0.14 & 11.1 & +0.298 & 23.7 & 0.29 & 13.4 & +0.270 & 14.6 & 0.12 & 10.9 & +0.295 & 27.0 & 0.29 \\
149 & 14.6 & +0.226 & 14.0 & 0.13 & 11.0 & +0.298 & 23.4 & 0.29 & 13.4 & +0.270 & 14.3 & 0.12 & 10.8 & +0.295 & 27.1 & 0.29 \\
\hline
\end{tabular}
\end{table}

\subsection{Later matched CPC V1--V2 DA ablation: a better background, but
similar analyses}
\label{sec:ablation-v1-v2-da}

The subsequent real-BMD experiment provides an important complement to the
training diagnostics above. To avoid version ambiguity, ``CPC V1'' in this
subsection denotes the operational checkpoint using \texttt{log1p}
precipitation transforms, and ``CPC V2'' denotes the later checkpoint using
square-root transforms; neither label refers to arms A1--A4. The comparison
uses 2022-05-01 through 2022-05-10, 38 BMD stations, five exhaustive spatial
folds, and 30 members. Each station is withheld exactly once. V1 and V2 are
aligned by station ID and date, and the uncertainty intervals use 10,000
paired circular-bootstrap resamples in three-day blocks, keeping all stations
from one weather day together.

This is a matched \emph{checkpoint-package} comparison, not a transform-only
experiment. Relative to V1, V2 changes the CHIRPS and CPC precipitation
transforms from \texttt{log1p} to \texttt{sqrt}, but also adds multiscale CPC
conditioning, weak coarse-consistency losses, wet-day/crop oversampling, and a
broader validation selection. The result is therefore consistent with a
benefit from the square-root formulation, but the gain cannot be attributed to
that transform alone.

\begin{table}[htbp]
\centering
\caption{Matched CPC V1--V2 performance at withheld BMD gauges. CRPS and MAE
are in mm~day$^{-1}$. V2 uses its selected spread-6, $\gamma=0.01$ gauge
guidance, while V1 retains the previously reported DA configuration. The
simultaneous rows use gauges plus the same S04 IMERG observation scale.}
\label{tab:ablation-v1-v2-da}
\small
\begin{tabular}{lrrrrrr}
\hline
Method & CRPS & MAE dry & MAE wet & Bias & Corr. & Cov.$_{90}$ \\
\hline
V1 background          & 7.547 & 10.38 & 17.63 & $+8.37$ & 0.362 & 0.81 \\
V2 background          & 5.915 &  7.01 & 13.73 & $+2.87$ & 0.400 & 0.94 \\
\hline
V1 gauges              & 4.961 &  4.10 & 11.86 & $+0.62$ & 0.537 & 0.58 \\
V2 gauges              & 4.872 &  4.12 & 11.76 & $+0.56$ & 0.540 & 0.77 \\
\hline
V1 simultaneous S04    & 4.768 &  4.10 & 12.19 & $+2.16$ & 0.565 & 0.64 \\
V2 simultaneous S04    & \textbf{4.626} & 2.90 & 11.86 & $+0.47$ & 0.569 & 0.84 \\
\hline
\end{tabular}
\end{table}

\begin{table}[htbp]
\centering
\caption{Paired fair-CRPS contrasts from the matched CPC V1--V2 experiment.
$\Delta=\mathrm{CRPS}(\mathrm{reference})-
\mathrm{CRPS}(\mathrm{candidate})$, so positive values favour the candidate.
Intervals crossing zero are reported as unresolved rather than wins.}
\label{tab:ablation-v1-v2-paired}
\small
\begin{tabular}{lrrl}
\hline
Candidate versus reference & $\Delta$ & 95\% interval & Decision \\
\hline
V2 background versus V1 background & $+1.632$ & $[+0.875,+2.476]$ & V2 better \\
V2 gauges versus V1 gauges & $+0.089$ & $[-0.108,+0.280]$ & unresolved \\
V2 simultaneous versus V1 simultaneous & $+0.142$ & $[-0.015,+0.301]$ & unresolved \\
V2 simultaneous versus V1 gauges & $+0.335$ & $[-0.057,+0.673]$ & unresolved \\
V2 gauges versus V2 background & $+1.043$ & $[+0.637,+1.542]$ & DA better \\
V2 simultaneous versus V2 gauges & $+0.246$ & $[-0.012,+0.457]$ & unresolved \\
\hline
\end{tabular}
\end{table}

\paragraph{The dominant V1--V2 difference is in the unassimilated prior.}
V2 reduces background CRPS by $1.632$~mm~day$^{-1}$, or $21.6\%$ relative to
V1, with a paired interval wholly above zero. The improvement is not confined
to one score: background bias falls from $+8.37$ to $+2.87$~mm~day$^{-1}$,
dry MAE from $10.38$ to $7.01$, wet MAE from $17.63$ to $13.73$, and
correlation rises from $0.362$ to $0.400$. Thus the V2 training package
materially improves the state \emph{before} any BMD or IMERG observation is
assimilated. This is the cleanest evidence that the prior redesign worked.

\paragraph{Assimilation compresses the checkpoint gap.}
After gauge assimilation, the V1--V2 CRPS difference shrinks from $1.632$ to
$0.089$~mm~day$^{-1}$; after simultaneous gauge--IMERG assimilation it is
$0.142$~mm~day$^{-1}$. Both paired intervals include zero, while correlations
are nearly identical (0.537 versus 0.540 for gauges and 0.565 versus 0.569 for
simultaneous). The absolute DA increment is consequently much larger for the
weaker V1 prior: gauges reduce V1 CRPS by $2.586$~mm~day$^{-1}$ but V2 CRPS by
$1.043$; simultaneous assimilation reduces them by $2.779$ and $1.289$,
respectively. The observations and likelihood therefore pull two unequal
backgrounds toward a similar observation-constrained analysis.

That convergence should not be read as evidence that prior quality is
irrelevant. First, V2 reaches the same analysis quality with substantially less
correction from the background, reducing dependence on sparse observations.
Second, the V2 simultaneous point estimate is the best of the six methods and
retains meaningful secondary advantages even though its CRPS gain over V1 is
unresolved: bias is $+0.47$ rather than $+2.16$~mm~day$^{-1}$, dry MAE is
$2.90$ rather than $4.10$, and 90\% coverage is $0.84$ rather than $0.64$.
Third, the V2 background is the field available away from observation
influence, so its improvement matters directly for spatial completeness even
when scores at withheld gauges converge.

\paragraph{Interpretation for the transform ablation.}
The defensible conclusion is two-part: the square-root-based V2 \emph{package}
produces a decisively better background, whereas the present DA configurations
take V1 and V2 to statistically indistinguishable withheld-gauge CRPS. This
suggests that the main practical value of V2 is improved prior quality rather
than a new DA response. It does \emph{not} establish that \texttt{sqrt} alone
is superior to \texttt{log1p}; isolating that causal claim requires retraining
an otherwise identical V2 configuration with only the precipitation transform
changed, followed by the same matched-fold background and DA evaluation.


\subsection{Findings}
\label{sec:abl-findings}

\paragraph{F1. No arm satisfies the wet-frequency criterion at convergence, and
all four are worse than the baseline.} The promotion criterion was a
wet-frequency error within $\pm 0.05$. At epoch 149 the errors are $+0.226$
(A1), $+0.298$ (A2), $+0.270$ (A3) and $+0.295$ (A4), against V1's $+0.19$.
Taken at face value, the redesign not only failed to fix the defect it targeted
but made it worse in every configuration tried. A2, A3 and A4 rain on more than
$96\%$ of cells at convergence.

\paragraph{F2. The target is nonetheless reachable: A1 attained a wet-frequency
error of $+0.002$ at epoch 39, and further training destroyed it.} This is the
central result. A1 is the only arm combining the hurdle head with
$\texttt{dry\_weight}=2.0$, and it is the only arm that was ever dry: it begins
at $0.629$ (error $-0.067$), passes through $0.698$ (error $+0.002$) at epoch
39, and then drifts monotonically to $0.922$. The mechanism therefore works. The
failure at epoch 149 is a failure of the training trajectory, not of the
proposed remedy, and the promotion criterion as written --- evaluated only at
the final or best-CRPS checkpoint --- would have discarded a configuration that
met the target perfectly a quarter of the way through training.

\paragraph{F3. Ensemble spread grows monotonically with training in every arm,
and is the proximate cause of the drift.} Spread against epoch correlates at
$r = +0.99$ (A1), $+0.96$ (A2), $+0.99$ (A3) and $+0.88$ (A4); no arm plateaus
within 150 epochs. Section~\ref{sec:abl-spread} treats this in full. In A1, the only arm with room to move on wet fraction, spread and wet
fraction correlate at $r = +0.96$. The interpretation is mechanical: as the
predictive distribution broadens, more cells and members cross the
$1$~mm~day$^{-1}$ threshold, so the measured wet fraction rises even where the
central estimate does not. A2 and A3 are saturated near $0.99$ from their first
validation and so show no such correlation; there is no headroom for one.

Growing spread is not an artefact of the metric. It is corroborated
independently by the assimilation diagnostics, where the background is
over-dispersed relative to its own innovations: the consistency ratio
$\mathrm{var}(O-B) / (\sigma_b^2 + R)$ is $0.62$, with background spread
variance $0.741$ against an innovation variance of $0.572$.

\paragraph{F4. The absolute target costs spatial skill.} Both absolute-target
arms converge to spatial correlations of $0.13$ (A1) and $0.13$ (A3), while the
residual arm A2 holds $0.29$. Since A3 has no hurdle head and A1 does, and the
two are indistinguishable on this measure, the degradation is attributable to
the target parameterisation rather than to the hurdle head. The residual
formulation supplies the CPC spatial structure to the model directly; CPC
remains a conditioning channel under the absolute parameterisation, but the
network does not learn to exploit it to the same degree. On the evidence here,
$T(\mathrm{CHIRPS}) - T(\mathrm{CPC})$ should be retained.

\paragraph{F5. The hurdle head alone is insufficient (partly superseded by F8).} A2 carries the hurdle head at $\texttt{dry\_weight}=1.0$ and sits at
$0.997$ wet fraction from its first validation --- indistinguishable from A3,
which has no hurdle head at all. Only A1, which combines the head with
$\texttt{dry\_weight}=2.0$, produces a dry prior. The head supplies the
mechanism for representing the atom at zero, but without upweighting the dry
cells in the flow loss the gradient signal is too weak for it to be used.

The stronger reading originally drawn from this --- that $\texttt{dry\_weight}$
was the \emph{sufficient} ingredient --- was tested by A4 and is wrong. See F8.

\paragraph{F6. A confound is present in all three arms, and it has now been
tested and rejected as the explanation.} All three v3 configurations set
$\texttt{cond\_dropout}=0.1$, whereas V1 used $0.0$. This was introduced to
enable a later unconditional-prior ablation and is unrelated to the wet-bias
hypothesis. Training with the conditioning dropped on $10\%$ of batches teaches
the model an unconditional branch, and all validation samples were drawn at
guidance weight $w=1$, i.e.\ with no classifier-free guidance applied to recover
the conditional sharpness. Since this is the only change common to all three
arms and absent from the baseline, it was the leading candidate explanation for
the excess spread, the inflated wet frequency and the depressed spatial
correlation.

It was tested directly by resuming the A2 checkpoint and resampling at $w=2$,
changing nothing else (Table~\ref{tab:cfg-probe}). Guidance reduced ensemble
spread by $14\%$ but left the wet-frequency error unchanged at $0.293$ against
$0.298$, and degraded both CRPS and spatial correlation. The hypothesis is
therefore rejected: the near-saturated wet fraction is not an artefact of
excess spread, and returning $\texttt{cond\_dropout}$ to $0.0$ should not be
expected to correct it. A2 is intrinsically wet.

\begin{table}[htbp]
\centering
\caption{Guidance probe on the A2 checkpoint. The $w=2$ row is at epoch 155
because the run must reach a checkpoint boundary to validate; over the preceding
ten epochs spread drifted by ${\sim}0.05$~mm~day$^{-1}$ per epoch, two orders of
magnitude below the effect under test.}
\label{tab:cfg-probe}
\begin{tabular}{lrrrr}
\hline
Guidance & Spread & Wet-frac. error & CRPS & Corr. \\
\hline
$w=1$ (epoch 149) & 23.35 & $+0.298$ & 10.97 & 0.286 \\
$w=2$ (epoch 155) & 19.98 & $\phantom{+}0.293$ & 12.69 & 0.257 \\
\hline
\end{tabular}
\end{table}

This result also sharpens F3. Spread and wet frequency are strongly coupled
\emph{within} A1 ($r=+0.96$), the one arm with headroom, but the coupling is not
universal: in A2, which is saturated near $0.99$, a $14\%$ reduction in spread
moves the wet fraction not at all. Excess spread amplifies the wet bias where
one is developing; it does not create it.

\paragraph{F7. Removing \texttt{cond\_dropout} at training time does not stop
the spread growth, confirming F6's rejection from the other direction.} A4 is
the only arm trained with $\texttt{cond\_dropout}=0.0$. Its spread nonetheless
rises from $5.35$ to $27.11$~mm~day$^{-1}$ over 150 epochs --- the largest of any
arm --- with $r=+0.88$ against epoch. The guidance probe rejected the
unconditional branch as the explanation at sampling time; A4 rejects it at
training time. The monotonic growth of predictive spread is therefore a property
of the training objective or schedule, and is now the principal unexplained
result of this study.

In A4 the coupling between spread and wet frequency reaches $r=+0.974$, the
strongest observed, which is consistent with the mechanical reading in F3: a
broadening distribution pushes more cells across the $1$~mm~day$^{-1}$ threshold.

\paragraph{F8. Correction to F5: the dry behaviour required the absolute target,
not the dry weighting alone.} F5 concluded from A1--A3 that
$\texttt{dry\_weight}=2.0$ was the active ingredient, since only the arm
carrying it was ever dry. A4 isolates that claim directly: it pairs the hurdle
head with $\texttt{dry\_weight}=2.0$ on the \emph{residual} target, and never
descends below a wet-frequency error of $+0.196$, against A1's $+0.002$. At
comparable spread ($5.35$ for A4 at epoch 9, $4.53$ for A1 at epoch 24) A1 is
substantially drier. The dry weighting is necessary but not sufficient; the
absolute parameterisation supplied the rest.

This leaves a genuine trade-off rather than a single best configuration. The
absolute target produces the correct wet frequency but poor spatial structure
(A1: correlation $0.13$); the residual target preserves structure (A4:
$0.29$--$0.40$, the best of the four) but is persistently too wet. A4 also
achieves the lowest CRPS of any arm ($9.38$ at epoch 39) while having the worst
wet-frequency error, which is a compact statement of why CRPS cannot be the
selection criterion for this defect.

\subsection{The spread pathology}
\label{sec:abl-spread}

Every arm follows the same trajectory: it begins comparatively dry and
under-dispersed, and both quantities rise together for the whole of training
without plateauing. Table~\ref{tab:ablation-spread} quantifies it.

\begin{table}[htbp]
\centering
\caption{Ensemble spread growth and its coupling to wet frequency. Spread is in
mm~day$^{-1}$; correlations are across the validation epochs of each run. No arm
plateaus within 150 epochs, and the arms with the most headroom on wet frequency
show the strongest coupling.}
\label{tab:ablation-spread}
\begin{tabular}{lrrrrr}
\hline
Arm & Spread first & Spread final & Growth & $r$(spread, epoch) & $r$(spread, wet frac.) \\
\hline
A1 & 4.53 & 14.05 & 3.1$\times$ & $+0.987$ & $+0.962$ \\
A2 & 14.89 & 23.35 & 1.6$\times$ & $+0.957$ & $-0.779$ \\
A3 & 2.24 & 14.29 & 6.4$\times$ & $+0.988$ & $-0.191$ \\
A4 & 5.35 & 27.11 & 5.1$\times$ & $+0.879$ & $+0.974$ \\
\hline
\end{tabular}
\end{table}

The negative correlations for A2 and A3 are not counter-evidence. Both are
saturated above $0.99$ wet fraction from their first validation, so the wet
frequency has nowhere to move and the residual variation is noise; the coupling
is only measurable where there is headroom, and there it is $+0.96$ (A1) and
$+0.97$ (A4).

Three explanations have been eliminated. Classifier-free guidance at $w=2$
reduces spread by $14\%$ and does not change the wet frequency
(Table~\ref{tab:cfg-probe}), so it is not a sampling artefact.
$\texttt{cond\_dropout}=0.0$ in A4 does not prevent it, so it is not the learned
unconditional branch. It appears with both target parameterisations and with and
without the hurdle head, so it is none of the components under test.

What remains is the training objective or schedule itself: the noise schedule,
the weighting of the loss across the flow time $t$, or the interaction between
the hurdle mask and the flow objective. A cheap next diagnostic is to measure
predictive spread as a function of $t$ on an existing checkpoint, which
localises the growth in flow time without training anything.

\subsection{Consequences for assimilation}

The prior's wet-frequency error is not confined to training diagnostics. In the
May--June 2024 assimilation experiment (427 withheld station-days, 159 of them
wet, gauges as truth), the rank histogram of the analysis ensemble carries
$0.62$ of its mass in the lowest rank --- the observation falls below every
member. The dry-day fraction of that sample is $0.628$. The agreement to three
decimal places identifies the cause: the ensemble effectively never predicts
zero, so on every dry day the gauge lies beneath the entire distribution. This
is the same defect, measured against gauges rather than CHIRPS, and it is the
mechanism by which it degrades the analysis.

The assimilation nonetheless improves substantially on the background (CRPS
$6.22$ against $9.45$; mean absolute error $9.56$ against $16.33$; bias $-0.66$
against $+7.99$~mm~day$^{-1}$). What it does not do is beat the raw CPC field it
was conditioned on ($9.62$ mean absolute error, $16.34$ on wet days against the
analysis's $18.50$, correlation $0.79$ against $0.76$). Closing that gap is the
practical reason the wet-frequency defect matters.

\subsection{Limitations}

\begin{enumerate}
  \item \textbf{The CPC V1--V2 result is neither transform-isolated nor an
    independent confirmation.} V2 changes the precipitation transform together
    with multiscale conditioning, loss terms and sampling, and the same
    2022-05-01--10 window was used to select its DA configuration. The paired
    intervals therefore describe sampling uncertainty conditional on those
    selected configurations; they do not include model-selection uncertainty.
    A frozen multi-season comparison is required before generalising the
    $21.6\%$ background gain beyond this development window.
  \item \textbf{The confound in F6 is untrained but tested.} No arm was trained
    with $\texttt{cond\_dropout}=0.0$, so the redesign is not cleanly separated
    from that change at training time. The sampling-side probe in F6 rejects it
    as the explanation for the wet-frequency failure, but does not exclude a
    training-side effect on the learned velocity field itself.
  \item \textbf{A1 varies two factors.} The full redesign changes both the
    target parameterisation and \texttt{dry\_weight} relative to A2, so the
    $+0.002$ result at epoch 39 cannot be attributed to the dry weighting alone
    without a fourth arm.
  \item \textbf{Six validation days.} The wet fraction is computed on six July
    days chosen to span the rainfall distribution. This is adequate for
    detecting an effect of the magnitude reported but not for fine comparison
    between arms.
  \item \textbf{Two arms were restarted.} A1 and A2 were cancelled at epoch 20
    and resumed from checkpoint after two defects were repaired in the
    validation path (the conditioning was not broadcast over the ensemble in the
    dry-logit evaluation, and the flow loss did not split the two-channel output
    when no dry target was supplied, so the reported validation loss was
    meaningless for hurdle models). Training itself was unaffected; their
    histories therefore begin at epoch 24 rather than 9.
  \item \textbf{Validation is referenced to CHIRPS.} This is appropriate here,
    since CHIRPS is the training target, but it is not a measure of skill
    against ground truth; the gauge-referenced evaluation is reported
    separately.
\end{enumerate}

\subsection{Implications}

The ablation supports retaining the residual target (F4) and pairing the hurdle
head with $\texttt{dry\_weight}=2.0$ (F5), since the head is inert on its own.
The proposal to return $\texttt{cond\_dropout}$ to $0.0$ was motivated by F6 and
is no longer supported by it; the change remains worth making to remove a
gratuitous difference from the baseline, but it should not be expected to move
the wet frequency, and a run adopting it should not be read as a test of that
hypothesis.

What remains unexplained is the ensemble spread, which grows monotonically
throughout training in all four arms without plateauing. The guidance probe
rules it out as a sampling artefact and A4 rules out $\texttt{cond\_dropout}$ as
a training cause, so the remaining candidates are the noise schedule, the loss
weighting across the $t$ range, and the interaction between the hurdle mask and
the flow objective. This is the open question the next experiment should
address, and on present evidence no configuration of target, hurdle head or dry
weighting will satisfy the wet-frequency criterion until it is resolved: every
arm is dry early and drifts wet as the spread grows.

Because the wet-frequency target is achieved transiently and then lost (F2, F3),
model selection must track wet-frequency error rather than CRPS alone; the
training loop now writes a separate wet-frequency-selected checkpoint for this
reason. On the v3 evidence CRPS continues to improve while the property of
interest degrades, so a CRPS-selected checkpoint is not the right one.
